\section{Discussion}
\label{sec:discussion}
\subsection{What Is Established by the Gate Ablation?}
The principal new finding is not merely that a small supervised head can improve PromptDA: a conventional residual CNN can do so as well. Rather, under an identical training objective and residual input interface, introducing a learned spatial gate improves over setting $G\equiv1$. The gated, no-gate-regularization control beats its ungated counterpart in all three seeds on ARKitScenes and all three paired seed/subset evaluations on ScanNet++ for AbsRel, L1, RMSE, RGB-edge-conditioned BoundaryAbsRel, and SSIM. Relative mean AbsRel improvements are 0.70\% and 1.03\% on the two datasets. The repeated direction supports a reproducible \emph{within-architecture} effect under the tested boundary-weighted objective. Nevertheless, the size of the improvements is modest and no scene-level uncertainty analysis or independent final test supports a statistical generalization claim.

An essential distinction is that this architectural attribution uses gate regularization disabled \emph{in both arms}, with the original boundary-weighted loss present. It is not an isolation of the gate in the selected no-boundary-loss SACG-NoB configuration. We avoid importing the result across differing objectives without a new matched experiment. The ungated branch also has fewer trainable parameters, so the comparison tests the benefit of the explicit gate architecture, not parameter efficiency or a fixed-parameter-budget superiority claim.

\subsection{Why a Selective Residual Is Plausible}
The residual branch can generate positive or negative metric corrections, while the gate determines their pixelwise magnitude using coarse depth, RGB--depth gradient-magnitude agreement, and LiDAR observation locations. This separation is conceptually useful when a pretrained model is already accurate in many regions. But the gate is not calibrated confidence: edge agreement can be high over textureless errors or low at purely photometric boundaries. The learned validity field identifies the spatial pattern of available LiDAR samples, not their measurement accuracy. Therefore, improvements in aggregate metrics do not prove that gate values correlate with true per-pixel error, or that the method specifically corrects errors caused by sensor-resolution mismatch.

Further measurements should stratify errors by distance to observed LiDAR, agreement score, actual ground-truth depth discontinuities, and the value of the learned gate. A particularly informative diagnostic would compare changes in both erroneous and already-correct pixels, since selective refinement should ideally avoid damaging the latter.

\subsection{Comparison With a Conventional CNN}
The expanded three-seed results prevent us from claiming SACG is globally better than a simple CNN refiner. On ARKitScenes, the CNN has slightly better mean AbsRel and L1 and lower RMSE and RGB-edge-based error than SACG-NoB, while its SSIM is also slightly higher. On ScanNet++, SACG-NoB has marginally lower mean AbsRel and RGB-edge-based error; the CNN is better on mean L1, RMSE, and SSIM. These patterns reinforce that performance is dataset- and metric-dependent.

Crucially, SACG and the separate CNN do not receive identical inputs. SACG uses frozen intermediate DPT features in its correction branch and a LiDAR-validity-derived field in its gate; the CNN receives resized numerical sparse LiDAR and a validity mask. Thus their performance differences cannot isolate the learned gate's utility. The direct constant-gate experiment resolves this attribution question within SACG's architecture, but further input-matched and parameter-matched CNN controls are still required to claim a general architectural advantage.

\subsection{Training Losses and Sensor Guidance}
SACG-NoB achieves better mean validation AbsRel and L1 than full SACG with the exponential RGB-edge-weighted boundary loss on the two evaluated datasets. The observation says only that this auxiliary loss did not help the chosen metrics under the recorded conditions, not that depth boundaries should be ignored. Similarly, learnable LiDAR-validity fields outperform fixed or removed fields in the full-objective ablations on mean AbsRel, but those conditions were not rerun under SACG-NoB. This prevents attribution of the selected configuration's advantage to a particular field parameterization.

The approach preserves PromptDA's original parameter values and changes only the appended adapter during training. That property makes it an attractive form of supervised downstream adaptation, but runtime, memory, and trainable parameter counts have not yet been quantified, and no genuine PromptDA fine-tuning baseline has been evaluated. Efficiency and superiority to full or partial fine-tuning therefore remain hypotheses.

\subsection{Limitations and Threats to Validity}
All reported comparisons use stored validation artifacts. The validation protocol also participated in checkpoint selection and the choice of objective, so the numbers must not be described as independent held-out test performance. ARKitScenes compares all seeds on identical 5,599 frames. In contrast, the ScanNet++ evaluation data seed changes with the model seed, meaning that its three-seed variation combines optimization and subset selection. The local ScanNet++ export also is not an official published benchmark replication.

BoundaryAbsRel is conditioned on RGB Sobel edges, which can reflect texture or lighting rather than actual geometric depth changes. The auxiliary SSIM uses the repository's invalid-depth handling, and neither metric establishes precise geometric edge preservation. Robust claims would require GT-depth discontinuity analysis, paired scene-bootstrap uncertainty, a final independent scene-held-out split, and input-matched controls. The current evaluation does not measure performance under changed sensor noise, sparsity, new devices, or real-world deployment conditions.

\section{Conclusion}
We examined Spatially Adaptive Consistency-Gated Refinement (SACG), a trainable residual adapter attached to a frozen PromptDA metric-depth model. Its correction branch uses RGB, coarse depth, and frozen decoder features, while a separate spatial gate uses image--depth gradient agreement, coarse depth, and LiDAR observation locations. Across three seeds on ARKitScenes and a local ScanNet++ export, the selected no-boundary-loss adapter reduces mean AbsRel and L1 relative to frozen PromptDA.

The main architectural result is a direct, loss-matched gate-network ablation: with boundary supervision retained and gate regularization disabled in both models, the learned gate improves over the same residual architecture with a constant gate in all six paired evaluations. Mean AbsRel reductions relative to the ungated control are 0.70\% on ARKitScenes and 1.03\% on ScanNet++. These findings support the value of a spatially modulated residual \emph{within the tested architecture}, without demonstrating broad superiority to CNN refinement: a simple CNN performs slightly better on mean ARKitScenes metrics. Independent held-out testing and mechanism-focused spatial error analysis are required before making stronger claims about robustness, geometric boundary preservation, or adaptation to low-cost LiDAR uncertainties.
